\subsection{ALTERNATIVE CAYLEY ALGEBRAS}

PROPOSITION 2. Let A be a ring, F a Cayley A-algebra, e its unit element, \(s: x \mapsto \bar{x}\) its conjugation and N: F \(\to\) A its Cayley norm (§ 3, no. 4).

\begin{enumerate}
\def\labelenumi{(\roman{enumi})}
\item
  For \(\mathbf{F}\) to be alternative, it is necessary and sufficient that, for every ordered pair of elements \(x, y\) of \(\mathbf{F}\) , \(x^2y = x(xy)\) .
\item
  If \(\mathbf{F}\) is alternative, then \(\mathbf{N}(xy) = \mathbf{N}(x)\mathbf{N}(y)\) for all \(x,y\) in \(\mathbf{F}\) .
\item
  Suppose that \(\mathbf{F}\) is alternative. For an element \(x \in \mathbf{F}\) to be invertible, it is necessary and sufficient that \(\mathbf{N}(x)\) be invertible in \(\mathbf{Ae}\) ; the inverse of \(x\) is then unique and equal to \(\mathbf{N}(x)^{-1}\bar{x}\) ; denoting this by \(x^{-1}\) ,
\end{enumerate}

\[
x ^ {- 1} (x y) = x \left(x ^ {- 1} y\right) = y
\]

or all \(y \in \mathbf{F}\) .

The condition \(x^{2}y = x(xy)\) is obviously necessary for F to be alternative (no. 1, Proposition 1). Conversely, if it holds for every ordered pair of elements of F, applying it to \(\bar{x}\) and \(\bar{y}\) , it gives \(\bar{x}^{2}\bar{y} = \bar{x}(\bar{x}\bar{y})\) ; applying the conjugation s to this relation, we obtain \(yx^{2} = (yx)x\) , so that conditions (c) of Proposition 1 of no. 1 are satisfied.

Obviously \(a(e, x, y) = 0\) for all \(x, y\) in \(\mathbf{F}\) . If \(\mathbf{F}\) is alternative, we therefore deduce from Proposition 1 (no. 1) that the subalgebra \(\mathbf{G}\) of \(\mathbf{F}\) generated by \(e, x\) and \(y\) is associative. As \(\bar{x} = -x + \mathrm{T}(x) \in -x + \mathrm{Ae}, \bar{x} \in \mathbf{G}\) and similarly \(\bar{y} \in \mathbf{G}\) . Then \(\mathrm{N}(xy) = (xy)(\overline{xy}) = xy.\bar{y}.\bar{x} = \mathrm{N}(y)x\bar{x} = \mathrm{N}(y)\mathrm{N}(x)\) , using the fact that \(\mathrm{N}(y) \in \mathrm{Ae}\) . This proves (ii).

Finally we prove (iii). If \(\mathrm{N}(x)\) is invertible in Ae and we write \(x' = \mathrm{N}(x)^{-1}\bar{x}\) , then \(xx' = x'x = e\) , for \(\mathrm{N}(x) = x\bar{x} = \bar{x}x\) . Conversely if x admits a left inverse \(x''\) , then \(\mathrm{N}(x'')\mathrm{N}(x) = \mathrm{N}(e) = e\) by (ii) and \(\mathrm{N}(x)\) is invertible in Ae; further, as \(x' = \mathrm{N}(x)^{-1}\bar{x}\) is in the subalgebra generated by x and e, the elements \(x, x', x''\) belong to the associative subalgebra generated by x, \(x''\) and e and hence \(x'' = x''(xx') = (x''x)x' = x'\) , whence the uniqueness assertion. The formulae \(x^{-1}(xy) = x(x^{-1}y) = y\) follow from the fact that \(x^{-1}\) , x and y are elements of the subalgebra generated by x, y and e, which is associative.

PROPOSITION 3. Let E be a Cayley A-algebra, \(\gamma\) an element of A and F the Cayley extension of E defined by \(\gamma\) and the conjugation of E (§ 2, no. 5, Proposition 5). For F to be alternative, it is necessary and sufficient that E be associative.

Let \(u = (x, y)\) , \(v = (x', y')\) be two elements of \(\mathbf{F}\) (where \(x, y, x', y'\) are in \(\mathbf{E}\) ). Then (§ 2, no. 5, formula (27))

\[
\left\{ \begin{array}{l} u ^ {2} v = ((x ^ {2} + \gamma \bar {y} y) x ^ {\prime} + \gamma \bar {y} ^ {\prime} (y \bar {x} + y x), (y \bar {x} + y x) \bar {x} ^ {\prime} + y ^ {\prime} (x ^ {2} + \gamma \bar {y} y)) \\ u (u v) = (x (x x ^ {\prime} + \gamma \bar {y} ^ {\prime} y) + \gamma (x ^ {\prime} \bar {y} + \bar {x}. \bar {y} ^ {\prime}) y, y (\bar {x} ^ {\prime} \bar {x} + \gamma \bar {y} y ^ {\prime}) + (y \bar {x} ^ {\prime} + y ^ {\prime} x) x). \end{array} \right.\tag{3}
\]

Using the fact that \(\bar{y}y\) and \(\bar{x} + x\) are in Ae, examining these formulae shows that the associativity of E implies \(u^{2}v = u(uv)\) and hence the fact that F is alternative (Proposition 2). Conversely, if F is alternative, the equation \(u^{2}v = u(uv)\) applied when \(y' = 0\) gives

\[
(y \bar {x} + y x) \bar {x} ^ {\prime} = y (\bar {x} ^ {\prime} \bar {x}) + (y \bar {x} ^ {\prime}) x.
\]

Now the left hand side is equal to \((y\mathrm{T}(x))\bar{x}' = y(\bar{x}'\mathrm{T}(x)) = y(\bar{x}'x + \bar{x}'\bar{x})\) ;

comparing with the right hand side, we obtain \((y\bar{x}^{\prime})x = y(\bar{x}^{\prime}x)\) , which proves the associativity of E, since x, y and \(\bar{x}^{\prime}\) are arbitrary elements of E.

